% ============================================================================
% LANGENTWURF LE-07c: Querer y Poder - Modalverben beim Einkaufen
% Spanisch Klasse 7 - Sequenz 7: Compras
% ============================================================================
% Tier: 1 (vollständig)
% Doppelstunde: 5-6 von 8
% Fachfarbe: #c41e3a (Karminrot)
% ============================================================================

\documentclass[11pt,a4paper]{article}

\usepackage[utf8]{inputenc}
\usepackage[T1]{fontenc}
\usepackage[ngerman]{babel}
\usepackage{geometry}
\usepackage{fancyhdr}
\usepackage{lastpage}
\usepackage{xcolor}
\usepackage{tcolorbox}
\usepackage{enumitem}
\usepackage{tabularx}
\usepackage{longtable}
\usepackage{array}
\usepackage{booktabs}
\usepackage{amssymb}

% ============================================================================
% KONFIGURATION
% ============================================================================

\newcommand{\fach}{Spanisch}
\newcommand{\fachfarbe}{c41e3a}
\newcommand{\jahrgangsstufe}{7}
\newcommand{\schulart}{Gesamtschule}
\newcommand{\stundenthema}{Querer y Poder -- Modalverben beim Einkaufen}
\newcommand{\reihenthema}{De compras}
\newcommand{\stundennummer}{7c}
\newcommand{\reihenstunden}{4}
\newcommand{\gerniveau}{A1}
\newcommand{\lehrwerk}{Qué pasa 1}
\newcommand{\lehrwerkseite}{S. 92-94}
\newcommand{\lerngruppengroesse}{24}

% ============================================================================
% LAYOUT
% ============================================================================
\geometry{left=2.5cm, right=2.5cm, top=2.5cm, bottom=2.5cm}
\definecolor{fach}{HTML}{\fachfarbe}
\definecolor{gray}{HTML}{666666}
\definecolor{lightgray}{HTML}{f5f5f5}
\definecolor{successgreen}{HTML}{2a7a4b}
\definecolor{warningorange}{HTML}{b35c00}
\definecolor{basis}{HTML}{2a7a4b}
\definecolor{standard}{HTML}{2c5aa0}
\definecolor{erweiterung}{HTML}{7b1fa2}

\tcbuselibrary{skins,breakable}

\newtcolorbox{infobox}[1][]{
    colback=lightgray,
    colframe=fach,
    fonttitle=\bfseries,
    title=#1,
    breakable
}

\newtcolorbox{scorebox}{
    colback=white,
    colframe=fach,
    fonttitle=\bfseries,
    title=Didaktik-Score,
    breakable
}

\pagestyle{fancy}
\fancyhf{}
\fancyhead[L]{\small\textcolor{gray}{\fach{} -- Jg. \jahrgangsstufe{} -- \stundenthema}}
\fancyhead[R]{\small\textcolor{gray}{LE-\stundennummer{} | Tier 1}}
\fancyfoot[C]{\small\thepage{} / \pageref{LastPage}}
\renewcommand{\headrulewidth}{0.4pt}

% Differenzierungsmarker
\newcommand{\B}{\textcolor{basis}{\textbf{[B]}}}
\newcommand{\Std}{\textcolor{standard}{\textbf{[S]}}}
\newcommand{\E}{\textcolor{erweiterung}{\textbf{[E]}}}

% ============================================================================
% DOKUMENT
% ============================================================================
\begin{document}

% ============================================================================
% DECKBLATT
% ============================================================================
\begin{titlepage}
    \centering
    \vspace*{2cm}
    {\large\textcolor{gray}{\schulart}}\\[2cm]
    {\Huge\textcolor{fach}{\textbf{Langentwurf}}}\\[0.5cm]
    {\large LE-\stundennummer{} | Tier 1 (vollständig)}\\[2cm]
    {\LARGE\textbf{\stundenthema}}\\[0.5cm]
    {\large Sequenz: \reihenthema}\\[2cm]
    \begin{tabular}{rl}
        \textbf{Fach:} & \fach{} (\gerniveau) \\[0.3cm]
        \textbf{Klasse:} & \jahrgangsstufe \\[0.3cm]
        \textbf{Lehrwerk:} & \lehrwerk, \lehrwerkseite \\[0.3cm]
        \textbf{Doppelstunde:} & \stundennummer{} (Std. 5-6 von 8) \\[0.3cm]
    \end{tabular}
    \vfill
\end{titlepage}

\tableofcontents
\newpage

% ============================================================================
% 1. BEDINGUNGSANALYSE
% ============================================================================
\section{Bedingungsanalyse}

\subsection{Lerngruppe}
\begin{tabular}{ll}
    Kursgröße: & \lerngruppengroesse{} SuS \\
    Jahrgangsstufe: & \jahrgangsstufe \\
    Sprachniveau: & \gerniveau{} (GER) -- fortgeschrittene Anfänger \\
    Fremdsprache: & 2. Fremdsprache (nach Englisch) \\
\end{tabular}

\subsection{Vorwissen aus der Sequenz}
\begin{itemize}
    \item \textbf{7a:} Einkaufsorte, Geschäfte (tienda, supermercado, mercado)
    \item \textbf{7b:} Kleidung, Farben (ropa, camiseta, pantalón, talla)
    \item Zahlen 1-100 (für Preisangaben)
    \item Regelmäßige Verben (-ar, -er, -ir)
    \item Verb tener (yo tengo, tú tienes...)
\end{itemize}

\subsection{Differenzierung (intern)}
\begin{tabular}{|l|p{3.5cm}|p{3.5cm}|p{3.5cm}|}
\hline
& \B{} \textbf{Basis} & \Std{} \textbf{Standard} & \E{} \textbf{Erweiterung} \\
\hline
\textbf{Ziel} & BR: Rezeption & MR: Produktion mit Hilfe & GY: Freie Produktion \\
\hline
\textbf{Konjugation} & Lückentext mit Vorgabe & Tabelle selbst ergänzen & + Stammwechsel erklären \\
\hline
\textbf{Dialog} & Nachsprechen & Mit Stichworten & Frei variieren \\
\hline
\end{tabular}

\vspace{0.5em}
\textbf{WICHTIG:} Diese Marker sind NUR für die Lehrkraft!

% ============================================================================
% 2. SACHANALYSE
% ============================================================================
\section{Sachanalyse}

\subsection{Sprachliche Strukturen}

\textbf{Verb \textit{querer} (wollen) -- Stammwechsel e $\rightarrow$ ie:}
\begin{center}
\begin{tabular}{|l|l|l|l|}
\hline
\textbf{Person} & \textbf{Pronomen} & \textbf{querer} & \textbf{Stammwechsel} \\
\hline
1. Sg. & yo & quiero & e $\rightarrow$ ie \\
2. Sg. & tú & quieres & e $\rightarrow$ ie \\
3. Sg. & él/ella/usted & quiere & e $\rightarrow$ ie \\
1. Pl. & nosotros/-as & queremos & (regelmäßig) \\
2. Pl. & vosotros/-as & queréis & (regelmäßig) \\
3. Pl. & ellos/ellas/ustedes & quieren & e $\rightarrow$ ie \\
\hline
\end{tabular}
\end{center}

\vspace{1em}

\textbf{Verb \textit{poder} (können) -- Stammwechsel o $\rightarrow$ ue:}
\begin{center}
\begin{tabular}{|l|l|l|l|}
\hline
\textbf{Person} & \textbf{Pronomen} & \textbf{poder} & \textbf{Stammwechsel} \\
\hline
1. Sg. & yo & puedo & o $\rightarrow$ ue \\
2. Sg. & tú & puedes & o $\rightarrow$ ue \\
3. Sg. & él/ella/usted & puede & o $\rightarrow$ ue \\
1. Pl. & nosotros/-as & podemos & (regelmäßig) \\
2. Pl. & vosotros/-as & podéis & (regelmäßig) \\
3. Pl. & ellos/ellas/ustedes & pueden & o $\rightarrow$ ue \\
\hline
\end{tabular}
\end{center}

\textbf{Redemittel -- Einkaufsdialog:}
\begin{itemize}
    \item \textbf{Verkäufer:} ¿Qué quiere(s)? / ¿En qué puedo ayudarle?
    \item \textbf{Kunde:} Quiero... / ¿Puedo ver...? / ¿Puedo probármelo?
    \item \textbf{Bezahlung:} ¿Puedo pagar con tarjeta? / ¿Cuánto cuesta?
\end{itemize}

\subsection{Kontrastive Betrachtung}

\begin{tabular}{|l|p{5cm}|p{5cm}|}
\hline
& \textbf{Deutsch} & \textbf{Spanisch} \\
\hline
Stammwechsel & Nicht vorhanden bei Modalverben & Systematisch: e$\rightarrow$ie, o$\rightarrow$ue \\
\hline
Modalverb + Infinitiv & ich will kaufen & quiero comprar \\
\hline
Höflichkeit & möchte (Konjunktiv) & quiero (Indikativ) ist normal \\
\hline
\end{tabular}

\textbf{Typische Interferenz:} *``Yo quero'' (ohne Stammwechsel) oder *``Yo puedo compro'' (finites Verb statt Infinitiv)

\subsection{Kulturelle Aspekte}
\begin{itemize}
    \item In Spanien ist ``Quiero...'' nicht unhöflich (anders als dt. ``Ich will'')
    \item Bezahlen mit Karte: ¿Se puede pagar con tarjeta? (unpersönliche Form üblich)
    \item Mercados in spanischsprachigen Ländern: Handeln ist teilweise üblich
\end{itemize}

% ============================================================================
% 3. DIDAKTISCHE ANALYSE
% ============================================================================
\section{Didaktische Analyse}

\subsection{Stellung in der Sequenz}
\begin{tabular}{|c|l|l|}
\hline
\textbf{Block} & \textbf{Thema} & \textbf{Status} \\
\hline
7a & Einkaufsorte + Geschäfte & abgeschlossen \\
7b & Kleidung + Farben & abgeschlossen \\
\textbf{7c} & \textbf{Querer y Poder -- Modalverben} & \textbf{diese Stunde} \\
7d & Compras-Test & folgt \\
\hline
\end{tabular}

\subsection{Kompetenzziele (Rahmenplan MV)}
\begin{itemize}
    \item \textbf{Kommunikativ:} Sprechen (an Gesprächen teilnehmen), Hörverstehen
    \item \textbf{Sprachlich:} Grammatik (Stammvokalwechsel), Wortschatz (Einkaufen)
    \item \textbf{Interkulturell:} Einkaufskultur in hispanischen Ländern
    \item \textbf{Methodisch:} Dialogisches Lernen, Rollenspiel
\end{itemize}

\subsection{Legitimation}
\textbf{Rahmenplan Spanisch MV (Kl. 7-10):}
\begin{itemize}
    \item Kompetenzbereich: Kommunikative Kompetenz -- Sprechen
    \item Standard: ``Die SuS können in einfachen Einkaufssituationen Wünsche äußern und Fragen stellen.''
    \item Grammatik: ``Die SuS können Stammvokalwechsel bei frequenten Verben anwenden.''
\end{itemize}

% ============================================================================
% 4. STUNDENZIELE
% ============================================================================
\section{Stundenziele}

\subsection{Grobziel}
Die SuS erweitern ihre \textbf{kommunikative Kompetenz im Bereich Sprechen}, indem sie Einkaufsdialoge führen und die Modalverben \textit{querer} und \textit{poder} mit Stammvokalwechsel korrekt anwenden.

\subsection{Feinziele (operationalisiert)}
\begin{tabular}{|c|p{7.5cm}|c|c|}
\hline
\textbf{Nr.} & \textbf{Die SuS können...} & \textbf{AFB} & \textbf{Diff.} \\
\hline
FZ1 & die Formen von \textit{querer} und \textit{poder} (yo, tú, él/ella) \textbf{nennen} und den Stammwechsel \textbf{erkennen} & I & alle \\
\hline
FZ2 & einen Einkaufsdialog mit vorgegebenen Redemitteln \textbf{führen} und Wünsche äußern & II & alle \\
\hline
FZ3 & die passende Form von \textit{querer/poder} in Kontexten \textbf{einsetzen} und die Wahl \textbf{begründen} & II & \Std{}/\E{} \\
\hline
FZ4 & einen eigenen Einkaufsdialog \textbf{entwickeln} und das Stammwechsel-Prinzip auf andere Verben \textbf{übertragen} & III & \E{} \\
\hline
\end{tabular}

\vspace{1em}
\textbf{AFB-Verteilung:} AFB I: 25\% | AFB II: 50\% | AFB III: 25\%

% ============================================================================
% 5. METHODISCHE ANALYSE
% ============================================================================
\section{Methodische Analyse}

\subsection{Einstieg}
\begin{itemize}
    \item \textbf{Methode:} Bildimpuls (Einkaufsszene) + Audio eines Einkaufsdialogs
    \item \textbf{Begründung:} Aktivierung des Vorwissens (7a/7b), Neugier auf neue Strukturen
    \item \textbf{Alternative:} Rollenspiel-Video (Verkäufer/Kunde)
\end{itemize}

\subsection{Erarbeitung}
\begin{itemize}
    \item \textbf{Methode:} Entdeckendes Lernen -- Stammwechsel aus Beispielen ableiten
    \item \textbf{Begründung:} Aktive Musterkennung statt passiver Rezeption
    \item \textbf{Scaffolding:} Farbcodierung der Stammvokale (qu\textcolor{fach}{e}rer $\rightarrow$ qu\textcolor{fach}{ie}ro)
\end{itemize}

\subsection{Übungsphasen}
\begin{itemize}
    \item \textbf{Methode:} Chorisches Sprechen $\rightarrow$ Partnerarbeit $\rightarrow$ Rollenspiel
    \item \textbf{Progression:} Gelenkt $\rightarrow$ Halbfrei $\rightarrow$ Frei
    \item \textbf{Differenzierung:} Redemittel-Karten mit/ohne Hilfen
\end{itemize}

\subsection{Medien}
\begin{itemize}
    \item Tafel/Whiteboard: Konjugationstabellen mit Farbmarkierung
    \item Lehrwerk: \lehrwerk, \lehrwerkseite
    \item Arbeitsblatt: AB-07c.pdf
    \item Lernpfad: LP-07c.html (Tablets, optional)
    \item Rollenkarten: Verkäufer/Kunde
    \item Realien: Kleidungsstücke oder Bilder für Rollenspiel
\end{itemize}

% ============================================================================
% 6. VERLAUFSPLANUNG
% ============================================================================
\section{Verlaufsplanung}

\textbf{1. Stunde (45 Min.)}

\begin{longtable}{|p{1cm}|p{1.5cm}|p{4cm}|p{3cm}|p{1.5cm}|p{2cm}|}
\hline
\textbf{Zeit} & \textbf{Phase} & \textbf{Lehrerhandeln} & \textbf{Schülerhandeln} & \textbf{SF} & \textbf{Material} \\
\hline
\endhead

0--5 & Einstieg & Bild: Tienda de ropa. ``¿Qué veis?'' Audio abspielen & Beschreiben, Vermutungen & Plenum & Bild, Audio \\
\hline

5--10 & Aktivier. & ``Was will die Person kaufen? Was hört ihr?'' & ``Quiero, puede...'' sammeln & Plenum & Tafel \\
\hline

10--20 & Input 1 & Verb \textit{querer}: Beispiele, Stammwechsel entdecken lassen & Muster erkennen, Tabelle ergänzen & Plenum & Tafel, AB \\
\hline

20--30 & Input 2 & Verb \textit{poder}: Analog zu querer einführen & Vergleichen, Tabelle ergänzen & Plenum & Tafel, AB \\
\hline

30--40 & Übung 1 & Chorisches Sprechen: ``Yo quiero, tú quieres...'' & Nachsprechen, Aussprache & Plenum & -- \\
\hline

40--45 & Sicherung & Merkkasten gemeinsam besprechen, Regel formulieren & AB-Merkkasten ergänzen & Plenum & AB-07c \\
\hline
\end{longtable}

\vspace{1em}

\textbf{2. Stunde (45 Min.)}

\begin{longtable}{|p{1cm}|p{1.5cm}|p{4cm}|p{3cm}|p{1.5cm}|p{2cm}|}
\hline
\textbf{Zeit} & \textbf{Phase} & \textbf{Lehrerhandeln} & \textbf{Schülerhandeln} & \textbf{SF} & \textbf{Material} \\
\hline
\endhead

0--5 & Aktivier. & Schnelles Quiz: ``Wie heißt \textit{ich will}?'' & Antworten: quiero & Plenum & -- \\
\hline

5--15 & Vertiefung & Lückentext AB, querer/poder einsetzen & AB bearbeiten & EA & AB-07c \\
\hline

15--25 & Anwendung & Rollenspiel: Verkäufer + Kunde, Einkaufsdialog & Dialoge mit Partnern üben & PA & Rollenkarten \\
\hline

25--35 & Transfer & \E{}: Weitere Stammwechselverben (pensar, dormir) & Muster übertragen & PA/GA & Tafel \\
\hline

35--40 & Sicherung & 2-3 Paare präsentieren Einkaufsdialog & Präsentieren, Feedback & Plenum & -- \\
\hline

40--45 & Abschluss & Zusammenfassung, HA: Vokabeln, LP-07c optional & Notieren & Plenum & -- \\
\hline
\end{longtable}

\textbf{Legende:} EA = Einzelarbeit, PA = Partnerarbeit, GA = Gruppenarbeit

% ============================================================================
% 7. ERWARTETE SCHWIERIGKEITEN
% ============================================================================
\section{Erwartete Schwierigkeiten}

\begin{tabular}{|p{4.5cm}|p{8cm}|}
\hline
\textbf{Schwierigkeit} & \textbf{Reaktion} \\
\hline
Stammwechsel vergessen: *quero & Farbmarkierung, Regel visualisieren (e$\rightarrow$ie bei betontem e) \\
\hline
Falscher Stammwechsel: *puiedo & Unterschied betonen: querer = e$\rightarrow$ie, poder = o$\rightarrow$ue \\
\hline
Infinitiv vergessen: *Quiero compro & Struktur verdeutlichen: Modalverb + Infinitiv! \\
\hline
nosotros/vosotros-Ausnahme & Merksatz: ``Wir und ihr -- bleiben fit!'' (kein Wechsel) \\
\hline
Zeitproblem & Puffer: Transfer-Aufgaben als HA aufgeben \\
\hline
\end{tabular}

% ============================================================================
% 8. HAUSAUFGABE
% ============================================================================
\section{Hausaufgabe}

\begin{itemize}
    \item Lehrwerk S. 93, Übung 4 (querer/poder einsetzen)
    \item Vokabeln lernen: Konjugation querer/poder, Einkaufs-Redemittel
    \item Optional: Lernpfad LP-07c.html bearbeiten
\end{itemize}

% ============================================================================
% 9. LÖSUNGEN
% ============================================================================
\section{Lösungen}

\subsection{AB-07c Lösungen}

\textbf{Merkkasten: querer}
\begin{itemize}
    \item yo $\rightarrow$ quiero
    \item tú $\rightarrow$ quieres
    \item él/ella $\rightarrow$ quiere
    \item nosotros $\rightarrow$ queremos
\end{itemize}

\textbf{Merkkasten: poder}
\begin{itemize}
    \item yo $\rightarrow$ puedo
    \item tú $\rightarrow$ puedes
    \item él/ella $\rightarrow$ puede
    \item nosotros $\rightarrow$ podemos
\end{itemize}

\textbf{Aufgabe 1: Konjugation zuordnen}
\begin{enumerate}[label=\alph*)]
    \item yo $\rightarrow$ quiero, puedo
    \item tú $\rightarrow$ quieres, puedes
    \item él/ella $\rightarrow$ quiere, puede
    \item nosotros $\rightarrow$ queremos, podemos
\end{enumerate}

\textbf{Aufgabe 2: Lückentext}
\begin{enumerate}[label=\alph*)]
    \item \textbf{Quiero} comprar una camiseta.
    \item ¿\textbf{Puedo} ver esa falda?
    \item Él \textbf{quiere} unos pantalones negros.
    \item Nosotros \textbf{podemos} pagar con tarjeta.
    \item ¿Tú \textbf{quieres} probarte los zapatos?
\end{enumerate}

\textbf{Aufgabe 3: Dialog (Beispiel)}
\begin{itemize}
    \item Vendedor: ¡Buenos días! ¿En qué puedo ayudarle?
    \item Cliente: Quiero una camiseta roja, por favor.
    \item Vendedor: ¿Qué talla quiere?
    \item Cliente: Talla M. ¿Puedo probármela?
    \item Vendedor: ¡Claro! El probador está allí.
    \item Cliente: Gracias. ¿Cuánto cuesta?
    \item Vendedor: Son 25 euros.
    \item Cliente: ¿Puedo pagar con tarjeta?
\end{itemize}

\textbf{Aufgabe 4: Stammwechsel erklären (Erweiterung)}
\begin{itemize}
    \item Regel: Bei betonter Silbe wechselt der Stammvokal
    \item e $\rightarrow$ ie: querer, pensar, entender
    \item o $\rightarrow$ ue: poder, dormir, volver
    \item Ausnahme: nosotros/vosotros (unbetont, kein Wechsel)
\end{itemize}

% ============================================================================
% 10. ANLAGEN
% ============================================================================
\section{Anlagen}

\begin{enumerate}
    \item Arbeitsblatt (AB-07c.pdf)
    \item Musterlösung (ML-07c.pdf)
    \item Lehrerhinweise (LH-07c.pdf)
    \item Lernpfad (LP-07c.html)
    \item Rollenkarten Vendedor/Cliente (Kopiervorlage)
    \item Bildkarten Kleidung (aus 7b)
\end{enumerate}

% ============================================================================
% 11. DIDAKTIK-SCORE
% ============================================================================
\newpage
\section{Didaktik-Score}

\begin{scorebox}
\textbf{Bewertung der didaktischen Qualität nach 10 Dimensionen}\\
\textbf{Ziel-Score: $\geq$ 9.0 (Premium-Qualität)}

\vspace{1em}
\begin{tabular}{|c|l|c|c|p{4cm}|}
\hline
\textbf{\#} & \textbf{Dimension} & \textbf{Gew.} & \textbf{Score} & \textbf{Notiz} \\
\hline
1 & Differenzierung & 15\% & 9/10 & [B]/[S]/[E] qualitativ \\
\hline
2 & Taxonomie (AFB) & 10\% & 9/10 & 25/50/25 erfüllt \\
\hline
3 & Lernziel-Alignment & 10\% & 10/10 & RP MV explizit \\
\hline
4 & Aktivierung & 10\% & 9/10 & Entdeckend, Rollenspiel \\
\hline
5 & Scaffolding & 10\% & 10/10 & Farbcodierung, gestuft \\
\hline
6 & Feedback-Qualität & 10\% & 9/10 & LP: elaboriert \\
\hline
7 & Sprachliche Klarheit & 10\% & 10/10 & Kl. 7 angemessen \\
\hline
8 & Motivation & 10\% & 9/10 & Alltagsrelevant (Einkaufen) \\
\hline
9 & Inklusion (UDL) & 10\% & 9/10 & Mehrere Zugänge \\
\hline
10 & Praktikabilität & 5\% & 9/10 & 90 Min. realistisch \\
\hline
\multicolumn{3}{|r|}{\textbf{GESAMT:}} & \textbf{9.3/10} & \\
\hline
\end{tabular}

\vspace{1em}
$\checkmark$ \textcolor{successgreen}{\textbf{FREIGABE} (Score $\geq$ 9.0)}
\end{scorebox}

\end{document}
